%%%%%%%%%%%%%%%%%%%%%%%%%%%%%%%%%%%%%%%%%%%%%%%%%%%%%%%%%%%%%%%%%%%%%%%%%%%%%%%%
%  SECTION: Context
%%%%%%%%%%%%%%%%%%%%%%%%%%%%%%%%%%%%%%%%%%%%%%%%%%%%%%%%%%%%%%%%%%%%%%%%%%%%%%%%
\section{Context}
\label{sec:Context}

The air transportation industry has been growing since 1945.
%
It passed through several economic and geopolitical crises, the most recent
being the COVID-19 pandemic~\cite{ACI_AirTravelDemand2025}.
%
The latter has significantly impacted the world passenger traffic, reducing the
industry by 40\%, when comparing 2021 against 2019, before the spread of the
virus~\cite{kalic2022impact}.
%
This diminish of the sector costs approximately US\$324 billion of gross
passenger operating revenue for airlines~\cite{ICAOEconomicImpact}.
% 
The COVID-19 also influenced the airport's revenue and operations.
%
At the beginning of the pandemic, the close sky policies allowed only essential
flights, reducing the use of the airport's infrastructure.
%
Sanitary measures also changed the aircraft movement itinerary, due to mandatory
cleaning and inspections before and after transporting passengers, which
directly affected the demand that could be accepted by the airports.
%



Even through all these negative indicators, it is estimated a recovery of almost
103\% of worldwide passenger transportation in 2024, when comparing against
2019~\cite{IATA2024}.
%
Some regions, such as Central America, have projections that surpass the
passenger demand of 2019 at the end of 2023.
%
Therefore, even with this uncertainty, the increasing demand required by the
airlines <<+ REFERENCE>>, to compose a better network, will result in congestion
of several airports, thereat old discussions of how to efficiently resolve
<<DETERMINE ?>> the distribution of aircraft operations at airports come back
into evidence.



To understand the operational concepts of an airport, it is essential to divide
it into two primary components: the airside and the landside~\cite{Lance2009}.
%
During an arrival, the airside process begins when air traffic management guides
an aircraft to land on the runway and taxi to its designated parking stand.
%
Once parked, passengers disembark and transition through the gates into the
landside terminal.
%
Inside the terminal, arriving passengers can utilize commercial facilities
before claiming their luggage and exiting via ground transportation.
%
Conversely, departing passengers begin at the landside by arriving via ground
transport, completing check-in and security, and moving through the terminal to
embark.
%
Finally, the aircraft returns to the airside by pushing back from the gate,
taxiing to the runway, and taking off.



With the prototypical airport conception, it is possible to start the discussion
about the slot allocation problem and how it is related to the way an airport
generate its revenue and the main limitations related to its infrastructure.



The airport revenue is mainly generated through aeronautical and
non-aeronautical sources, both of which are directly related to the aircraft
traffic volume.
%
The former is composed by the fare paid by each passenger using the airport
infrastructure, usually included in the flight cost paid to the airlines; and
the movement charge, which is paid by the airline when it lands, uses the
infrastructure (e.g., hangar and terminal), and takes-off from the airport.
%
The latter is the revenue generated by the commercial concession at the
terminal such as rental fees for the retail at the airport area; the real
estate exploitation, which includes hotels, restaurants, and parking lots at the
airport perimeter; and services offered to passengers and airlines.
%
A market equilibrium between both revenues sources is needed, since reduction
of the aeronautical charges commonly means stimulation to the non-aeronautical
revenues, through passenger demand increments, but it does not guarantee an
overall maximum return~\cite{ICAO2013}.



In 2019, the european aeronautical sources generated 54\% of overall airport
revenue~\cite{ACI2021}, with 63\% of that aeronautical income originating from
charges directly applied to passengers~\cite{ICAO2013}.
%
Furthermore, airport users often do not pay the full cost of the infrastructure
they utilize, as evidenced by the fact that 69\% of airports worldwide were
loss-making in 2012~\cite{ACI2021}.
%
To maximize profitability, airports have expanded their variety of commercial
activities, which cushions the impact of lower passenger and freight volumes due
to the higher profit margins associated with these services.
%
This creates a cross-subsidization between aeronautical and non-aeronautical
revenues sources, since profits from commercial exploitation, for example, can
be reinvested in airport infrastructure, increase the demand of passenger or
aircraft that can be accepted at the airport, reducing capital needs and overall
cost.



Knowing that airports are important for a country social-economic-growth and
they can achieve profitability with an efficient management, governments have
decided that, under the right economic conditions, they can include private
sector participation (e.g., outright ownership, short or long-term concessions,
public-private-partnership schemes, and management contracts) for the financing
and operation of an airport infrastructure.
%
The airports, under this new ownership model, have sought for better ways to
compete for both passenger and airlines.
%
The aviation community agrees with the need of infrastructure investments to
accommodate the growth of the industry, increasing the profit margins.
%
Since most of the infrastructure added are in large increments, and the time
needed for these ventures is extensive, the airports are exposed to considerable
risk~\cite{ICAO2013}.
%
Even with the constant pressure for the airport expansion, most of the
administrator have to better manage its current infrastructure, until the
investment risks are reduced to an acceptable level.
%
When the expansion of infrastructure is not possible, the optimization of the
allocation in the available slots is strongly sought by the airport
administrator, to improve the revenue generation prospects.
%
Some regions have already implemented more restrictive on-time performance goals
for airlines~\cite{ravizza2014aircraft}, which aims to enhance aircraft's taxi
times, take-off sequencing, and allocating push-back time, thus affecting the
demand of movements that an airport can accept.



The limitation of airport infrastructure use is determined by several
interconnected factors, including the number of aircraft movements runways can
accept within specific time intervals and the availability and category of
aircraft stands~\cite{ARC2021, WASG2020}.
%
Capacity is further defined by the number of gates available for passenger
movements, the processing limits of security inspections at check-in, and the
total passenger volume the terminal area can support.
%
Additionally, infrastructure use is constrained by the number of aircraft
permitted through time slots and specific environmental restrictions.
%
Ultimately, the most restrictive infrastructure capacity -- or a combination of
these various limitations -- is used to manage airport movements and distribute
slots effectively without compromising the level of service through delays or
taxiway congestion.


<<SLOTS APPEARS FIRST HERE. GIVE AN BRIEF DEFINITION>>






%%%%%%%%%%%%%%%%%%%%%%%%%%%%%%%%%%%%%%%%%%%%%%%%%%%%%%%%%%%%%%%%%%%%%%%%%%%%%%%%
%  SECTION: Problem Definition
%%%%%%%%%%%%%%%%%%%%%%%%%%%%%%%%%%%%%%%%%%%%%%%%%%%%%%%%%%%%%%%%%%%%%%%%%%%%%%%%
\section{Problem Definition}
\label{sec:ProblemDefinition}

The present study introduces two main reasons to raise the efficiency in the
initial slot allocation in an airport.
%
The first reason is related to the airport's ways to generate revenue.
%
The two main ways of an airport generates revenue (i.e., aeronautical and
non-aeronautical) are directly related to the passenger traffic volume.
%
Commercial revenues play a vital role in global airport economics, accounting
for $52.9\%$ of total revenue in Africa and the Middle East, $52.6\%$ in North
America, $45.7\%$ in Asia-Pacific, and $38.1\%$ in Europe, yet they comprise
only $29.0\%$ in Latin America and the Caribbean, where a lack of traffic
stimulus often results in a heavy dependence on regulated fares and
underdevelopment in both commercial and aeronautical
sectors~\cite{graham2009important}.
%
In the Economic Commission for America Latin and the
Caribbean~\cite{planzer2019airport}, it is demonstrated the growth of the
worldwide air transportation market before the COVID-19 pandemic (2017
<<REFERENCE?>>) which, by \acrfull{ICAO}, would reach 10 billion passengers per
year by 2040.
%
Despite the great projection, the report also describes the Latin American and
the Caribbean challenges to build airport infrastructure, the major impediment
to the development of the local sector.
%
The \acrfull{IATA} has urged the local aviation authorities to ensure, between
several highlighted topics, the efficient use of the available airport’s
capacity in the region.



The second reason is related to the impact of the operational performance
improvement by the airlines.
%
The operational capacity of an airport, i.e., the ability to absorb the demand
of flights, is directly related to its infrastructure (e.g., number of runways,
gates, aircraft stands, and terminal areas).
%
Although the mentioned features of an airport can be used to determine its
capacity, a lower limit is set to mitigate the negative impacts of the airport’s
main partners, the airlines.
%
This measure <<SINONYM?>> goes in the opposite direction of the first reason.
%
When the limit is set greater than it should be, operational impacts are raised
(e.g., delays).
%
Since most airlines interconnect their flights within a network, a single delay
can spread repercussions far beyond the local level, impacting multiple aircraft
itineraries and significantly elevating an airline’s operational costs.
%
This topic is discussed in~\cite{ball2010total}, where it is described that the
airline's costs of crew management, disrupted passenger accommodation and
aircraft re-position, due to delays, can reach 20 billion dollars a year (2017).
<<REWRITE LAST SENTENCE TO MOVE THE REFERENCE TO THE END OF THE SENTENCE.>>
%
Another negative impact is related to the intense and unplanned aircraft taxi
traffic.
%
Fuel consumed in stop-and-go situations, primarily resulting from congestion on
the airport’s taxiway system, can account for approximately 18\% of the fuel
used during the entire operation \cite{nikoleris2011detailed}.




Therefore, while both airports and air operators want to expand their
operations, a balance between accepted demand, which increases revenues (for
airports), and operational efficiency, which mitigates high operational costs
(for airlines), is needed.


<<+ PROBLEM STATEMENT: HOW CAN THE EFFICIENCY RAISE OF SLOTS ALLOCATION
CONTRIBUTES TO THE EFFICIENCY OF AIRPORT OPERATION AND ITS REVENUE GENERATION?>>














%%%%%%%%%%%%%%%%%%%%%%%%%%%%%%%%%%%%%%%%%%%%%%%%%%%%%%%%%%%%%%%%%%%%%%%%%%%%%%%%
%  SECTION: Justificative
%%%%%%%%%%%%%%%%%%%%%%%%%%%%%%%%%%%%%%%%%%%%%%%%%%%%%%%%%%%%%%%%%%%%%%%%%%%%%%%%
\section{Justificative}
\label{sec:Justificative}





%%%%%%%%%%%%%%%%%%%%%%%%%%%%%%%%%%%%%%%%%%%%%%%%%%%%%%%%%%%%%%%%%%%%%%%%%%%%%%%%
%  SECTION: Objectives
%%%%%%%%%%%%%%%%%%%%%%%%%%%%%%%%%%%%%%%%%%%%%%%%%%%%%%%%%%%%%%%%%%%%%%%%%%%%%%%%
\section{Objectives}
\label{sec:Objectives}





%%%%%%%%%%%%%%%%%%%%%%%%%%%%%%%%%%%%%%%%%%%%%%%%%%%%%%%%%%%%%%%%%%%%%%%%%%%%%%%%
%  SECTION: Text Structure
%%%%%%%%%%%%%%%%%%%%%%%%%%%%%%%%%%%%%%%%%%%%%%%%%%%%%%%%%%%%%%%%%%%%%%%%%%%%%%%%
\section{Text Structure}
\label{sec:Text Structure}





